\chapter{Kubernetes: Giải phẫu cụm và vòng lặp điều hòa}
\label{ch:cluster}

\section{Ý tưởng duy nhất mà mọi thứ khác treo lên}

Bạn không bảo Kubernetes \emph{làm} gì; bạn bảo nó điều gì phải
\emph{đúng}. Bạn ghi \emph{trạng thái mong muốn} (``một Deployment tên
course-service với 1 replica của image X'') vào cơ sở dữ liệu của nó
qua API server; các \emph{controller} chạy vòng lặp điều hòa
(reconciliation) mãi mãi, so sánh mong muốn với thực tế và hành động
trên phần chênh lệch. Xóa một pod: controller của nó nhận ra thực tế
$\neq$ mong muốn và tạo cái khác. Không ai phát lại lệnh của bạn ---
trạng thái tự hội tụ.

Điều này giải thích những hành vi nếu không sẽ có vẻ lạ:

\begin{itemize}
  \item \code{kubectl apply} là một phép gộp vào trạng thái mong muốn,
  không phải chạy script --- apply cùng một file hai lần, lần thứ hai
  chẳng có gì xảy ra (tính lũy đẳng --- thuộc tính mà pipeline dựa vào
  ở Chương~23).
  \item Không có ``thứ tự khởi động.'' Mọi thứ khởi động cùng lúc;
  service nào cần dependency chưa có sẽ crash và được khởi động lại cho
  tới khi thế giới sẵn sàng. Crash-loop \emph{trong lúc} hội tụ là vận
  hành bình thường, không phải lỗi.
  \item Xóa khỏi file manifest không xóa khỏi cụm --- apply chỉ thêm và
  cập nhật, nó không dọn (prune). Gỡ Ingress course-service cũ đã cần
  một lệnh \code{kubectl delete} tường minh.
\end{itemize}

\section{Các bộ phận chuyển động}

\begin{figure}[h]
\centering
\begin{tikzpicture}[node distance=4mm and 10mm]
  \node[boxdark] (api) {API server};
  \node[boxfill, right=of api] (etcd) {state store\\\scriptsize (etcd / sqlite)};
  \node[box, below left=6mm and -2mm of api] (sched) {scheduler};
  \node[box, below right=6mm and -2mm of api] (cm) {controller\\manager};
  \node[box, below=16mm of api] (kubelet) {kubelet + containerd\\\scriptsize trên node};
  \draw[flow] (api) -- (etcd);
  \draw[flowdash] (sched) -- (api);
  \draw[flowdash] (cm) -- node[lbl,right]{vòng điều hòa} (api);
  \draw[flow] (kubelet) -- node[lbl,left]{watch \& report} (api);
  \node[lbl, left=8mm of api] (kubectl) {kubectl};
  \draw[flow] (kubectl) -- (api);
\end{tikzpicture}
\caption{Mọi thứ chỉ nói chuyện với API server.}
\end{figure}

\begin{description}
  \item[API server] cửa trước duy nhất; \code{kubectl} là một REST
  client của nó. Xác thực, phân quyền (RBAC của Chương~22) và kiểm tra
  hợp lệ diễn ra ở đây.
  \item[Kho trạng thái (state store)] cơ sở dữ liệu trạng thái mong
  muốn (etcd trong cụm chuẩn; k3s nhúng sqlite khi chỉ có một node).
  \item[Bộ lập lịch (scheduler)] gán pod chưa được lập lịch vào node,
  tôn trọng \emph{requests} tài nguyên của Chương~17.
  \item[Controller manager] các vòng lặp điều hòa: Deployment
  $\rightarrow$ ReplicaSet $\rightarrow$ pod, endpoints, và hàng chục
  vòng khác.
  \item[kubelet] agent trên node: theo dõi các pod được gán cho node
  mình, điều khiển containerd để chạy chúng, thực thi probe, báo cáo
  trạng thái.
\end{description}

k3s đóng gói toàn bộ những thứ này vào một binary --- một Kubernetes
được CNCF chứng nhận, trừ đi chi phí của một VM, cộng thêm pin sẵn:
Traefik làm Ingress controller, lưu trữ local-path, containerd. Mọi
manifest trong sách này áp dụng nguyên xi lên EKS hay OpenShift; chính
chứng nhận đó là lý do học trên k3s chuyển giao được.

\section{Object, label, namespace}

Mọi object là một tài liệu YAML với cùng bộ khung --- \code{apiVersion},
\code{kind}, \code{metadata} (tên, namespace, label), \code{spec} (mong
muốn), \code{status} (quan sát được, do controller ghi). Hai cơ chế tổ
chức:

\begin{description}
  \item[Namespace] phân vùng cụm; dự án này đặt mọi thứ trong
  \code{ts}, giữ nó tách biệt với \code{kube-system} và cho phép giới
  hạn phạm vi RBAC --- quyền của Jenkins dừng lại ở ranh giới namespace
  (Chương~22).
  \item[Label] cặp khóa/giá trị tùy ý; \emph{selector} truy vấn chúng.
  Mọi mối nối trong Phần~V đều là một label selector: Service tìm pod
  theo label, Deployment sở hữu pod theo label. Kustomization đóng dấu
  \code{app.kubernetes.io/part-of: teacher-supporter} lên mọi thứ để
  một selector có thể liệt kê --- hoặc xóa --- cả ứng dụng.
\end{description}

\section{Đưa kubectl tới cụm}

\code{kubectl} đọc \code{\textasciitilde/.kube/config}: URL của cụm,
chứng chỉ CA và một thông tin đăng nhập. Từ Windows, dự án này với tới
\code{https://100.85.66.43:6443} qua Tailscale --- điều đòi hỏi cài k3s
với \code{--tls-san 100.85.66.43}, vì chứng chỉ của API server chỉ hợp
lệ cho những tên được nướng vào lúc cài đặt.

\begin{handson}[cảm nhận vòng lặp]
\begin{lstlisting}[language=bash]
kubectl -n ts get pods
kubectl -n ts delete pod -l app=course-service
kubectl -n ts get pods -w
\end{lstlisting}
Xóa pod, rồi quan sát: một pod thay thế xuất hiện trong vài giây, được
tạo không phải bởi lệnh của bạn mà bởi ReplicaSet controller nhận ra
thiếu một replica. Bạn chưa hề ra lệnh khởi động lại --- bạn làm trạng
thái thực tế lệch khỏi mong muốn, và vòng lặp đã khép lại khoảng cách
đó.
\end{handson}

\begin{pitfall}[certificate is valid for\ldots{} not 100.85.66.43]
Triệu chứng: \code{kubectl} từ Windows thất bại khi kiểm tra TLS trong
khi chạy tốt trên server. Nguyên nhân: chứng chỉ API server thiếu địa
chỉ tailnet trong SAN --- nó được cài mà không có \code{--tls-san}. Sửa
lúc cài đặt chỉ là một cờ; sửa sau đó nghĩa là xóa thư mục TLS của
server và khởi động lại k3s. Một trường hợp cụ thể của quy tắc chung:
chứng chỉ gắn với \emph{tên}, nên hãy quyết định mọi tên mà server sẽ
được gọi \emph{trước} khi cấp.
\end{pitfall}

\begin{quiz}
\begin{enumerate}
  \item Giải thích ``khai báo'': bạn viết gì, ai hành động trên nó, và
  vì sao apply cùng một manifest hai lần không thay đổi gì?
  \item Vì sao Kubernetes cố ý không có \code{depends\_on}, và điều gì
  thay cho thứ tự khởi động trong lúc hội tụ?
  \item Lần theo \code{kubectl -n ts delete pod ...} qua các thành
  phần: ai nhận, ai nhận ra, ai hành động, trên máy nào?
  \item Vì sao gỡ Ingress cũ khỏi git không gỡ nó khỏi cụm, và hai lệnh
  nào đều có thể gỡ được nó?
\end{enumerate}
\end{quiz}
